\clearpage
\section{Diagram Contracts and Reading Guide}
\label{sec:diagram-contracts}
Each diagram has one question and an explicit abstraction boundary. Solid arrows
show the operation or reference named on the edge; dashed arrows have a local
meaning given in the caption. A logical dependency does not imply an HTTP call,
independent service or physical foreign key. The sequences are numbered
collaboration summaries rather than complete UML models of every exception.

\subsection{Course Lifecycle: Who Can Make the Twin Available?}
Figure~\ref{fig:lifecycle} begins with a professor's owned course, permitted
materials and teaching settings. The output of configuration is a release draft,
not a student-visible chatbot. Preview lets the professor inspect behaviour;
preflight returns specific blockers. Only a subsequent publish operation makes
the release current. A student must independently pass membership and release
checks before opening a compatible conversation. Check-in consent is a separate
preference, not an automatic consequence of asking a question. Saved difficulty
signals and action outcomes inform the instructor's next review; they do not
change published policy automatically. Withdrawal blocks subsequent use of that
release. This figure establishes actor responsibilities; the publication
sequence supplies the transactional detail.

\subsection{Logical Architecture: Where Does Each Decision Execute?}
Figure~\ref{fig:responsibilities} separates web requests from worker entry points.
Routes authenticate and validate transport inputs; publication, tutoring and
autonomy services enforce domain rules. The ingestion worker processes queued
source jobs. The autonomy worker observes durable events and processes due
opportunities; it does not depend on an open browser tab. Both use application
services and repository contracts. The local deployment uses SQLite and source
files; LangGraph checkpoints support workflow recovery. Model adapters return
proposals, not authorization or trusted database updates. Provider failures
therefore enter service/graph recovery paths. This is a responsibility map:
separate boxes do not claim microservice deployment, independent scaling or
separate physical databases. A deployment diagram would answer a different
question and is not inferred from this figure.

\subsection{Response Sequence: What Must Hold Before a Reply Is Saved?}
Figure~\ref{fig:response-sequence} accepts a conversation identifier, content and
client request identifier. Scope checks resolve the student, membership and
release. Retrieval returns candidate chunks with source identity and locators;
the evidence gate selects what the generator may actually use. This boundary
explains why complete retrieval alone did not make the any-hit control answer.
The generator receives the selected evidence, policy and teaching intent. The
retained factual generator is deterministic; configured candidates may use an
external model. The returned text and claims remain proposals until claim and
citation checks finish.

The service then prepares messages, citations, audit information and applicable
learner and follow-up artifacts. The repository checks current authority and the
expected learner-state revision during the turn commit. A matching duplicate
request can return the saved response; reuse with different content is rejected.
A stale revision or withdrawn scope cannot silently overwrite current state.
Insufficient evidence leads to clarification or abstention, while policy-restricted
requests use refusal or redirection. Failed generation follows the configured
bounded recovery path. Model execution and checkpoint writes need not share the turn transaction.

\subsection{Autonomy Activity: What Happens Between Two Student Turns?}
Figure~\ref{fig:autonomy} starts with an observed event or scheduled wake-up.
The service creates or loads a scope-bound opportunity and claims it using an
expiring lease. The job snapshot includes goal status, permitted evidence,
release/profile identity, consent, quiet hours and recent-message constraints.
Terminal goals stop before further work. The planner intersects event actions
with instructor permissions, compares the deterministic fallback with a guarded
model proposal, and consumes the default delivery/event proxy unless an explicit
state resolver is supplied. The graph authorizes the resulting action,
generates wording and validates the result. A blocked or invalid proposal
produces a no-action/failure result; the worker is not obliged to send anything.
A valid response goes through the outreach service, which rechecks authority
before persisting an in-app message.

Delivery and final job-result commit are distinct stages. A stable trigger key
links retries to the same delivery instead of issuing a second message after
an interrupted attempt. The action outcome and subsequent learner observations
support later goal interpretation. The corrected V2 goal interpreter uses
the exact release-bound objective and requires every target concept to satisfy
the committed-evidence rule; the planning proxy cannot substitute for that
assessment. The dashed return requires a new durable
event or due wake-up: it is not unrestricted model self-invocation. Completion,
cancellation, expiry and attempt limits terminate work for that goal. The graph
explains operational control; its boxes do not establish that learner estimates
or teaching interventions are educationally effective.

\subsection{Publication Sequence: What Changes Atomically?}
Figure~\ref{fig:publication-detail} distinguishes preflight from publication.
Preflight checks active component-profile compatibility, approved policy,
evidence readiness, source permissions and unique provenance. Where configured,
index preparation also contributes to readiness. It persists the check outcome
and audit information. Its evaluation status is a software readiness field;
a passing status does not mean that all academic accuracy targets have passed.

Publication rechecks ownership and publishability, prepares the index and asks
the repository to replace the current release. The repository serializes the
transition, rechecks applicable teaching-profile authority, withdraws the
predecessor and cancels obsolete autonomous work. A post-publication observer
runs after that transition: failure there must not be interpreted as reversal
of a committed release. Rollback is revalidation and republication of a withdrawn
release, not restoration of an entire database. Index preparation, the database
transition and the observer hook are distinct failure boundaries.

\subsection{Data Relationships: Which Identity Survives a Change?}
Figure~\ref{fig:data-detail} uses explicit identifiers to explain persistence.
Membership associates an account with a course; a conversation references a
specific release; messages belong to conversations; citations retain source
identity, version, checksum and location. Release-bound policy and evidence
provide a snapshot against which an interaction can later be interpreted.
Learner observations and beliefs are derived records with scope and revision,
not ground truth copied from the source. Goals retain objectives and stopping
conditions; opportunities retain event windows and optional goal references;
actions/outcomes identify what processing actually did. The optional goal link
matters because an opportunity need not always originate from a saved goal.
Dashed update arrows describe information flow, not SQL referential constraints.
The figure deliberately omits unrelated fields and does not merge the grouped
plan/action/outcome records into one physical table.
